\section{Expected practical impact}
\label{sec:impact}

Nothing in this section has happened. \evhyp{} Every benefit below is conditional on a named study
reading continue, and none follows from the PoC alone: the PoC produced gap candidates, and no human
has checked one. \Cref{tab:impact} names who would benefit, what they would get, and what would have
to be true first.

\begin{table}[!htbp]
  \centering
  \footnotesize
  \caption[Expected practical impact]{Expected practical impact, by beneficiary. Each row holds only
    if the gate in the last column reads continue (\cref{sec:validation,sec:roadmap}).}
  \label{tab:impact}
  \begin{tabularx}{\linewidth}{@{}>{\raggedright\arraybackslash}p{0.2\linewidth}
      >{\raggedright\arraybackslash}X>{\raggedright\arraybackslash}p{0.3\linewidth}@{}}
    \toprule
    Who & What they would get & What must be true first \\
    \midrule
    CTA practitioners, knowledge engineers, training designers &
      A ranked list of likely gaps before the first interview, each with a predicted knowledge type
      and a suggested elicitation channel; less expert time per validated item &
      The gap map predicts on fresh, blind elicitation (V5, MS5), and targeting saves time by at
      least the proposed 1.25 times (V6, MS6) \\
    \addlinespace
    Software organizations &
      A knowledge-at-risk view per component that adds artifact coverage to authorship
      concentration, in aggregate form &
      Coverage predicts post-departure trouble beyond authorship (V3, the separate organizational
      hypothesis H-OSS), and a partner pilot is legal and beats controls (V7, MS7) \\
    \addlinespace
    Maintenance and compliance organizations &
      Targeted capture of diagnostic cues and exception judgments before a senior person leaves &
      The method holds in a text-poor diagnostic domain and in regulatory judgment (V5 in D3 and
      D4) \\
    \addlinespace
    Educators and learners &
      Instruction that states validated hidden steps, such as principle selection &
      Residual areas carry more novice errors (V8c, MS8), and a trial shows a delayed-transfer
      effect \\
    \addlinespace
    Researchers &
      Open closure set, reconstruct-versus-CTA benchmark, coded departure dataset &
      Only that V1, V2 and V3 run; each is released whatever its verdict \\
    \bottomrule
  \end{tabularx}
\end{table}

\paragraph{The size of any benefit is unknown, and some impact would remain even if H2 fails.}
\evint{} The saving per validated item is unknown until V6 measures it. Whether a coverage score
adds anything to authorship data is exactly what V3 tests. We found no credible, peer-reviewed cost
figure for organizational knowledge loss in general, so we claim none. Field effects of education
interventions are usually small \parencite{kraft2020interpreting}, so even a success brings modest
gains. Two fallbacks keep some value if H2 fails. If the closure study (MS1) stops the automated
absence check, the reconstruction pipeline still works as a documentation audit. If the screening
gate (MS3) does not release Phase~2, H2 funding stops and the separately costed interview-only pivot
follows; a positive result there would give practitioners an elicitation-efficiency tool, with no
claim about text (\cref{sec:risks}). Neither is a residual predictor, and neither would be sold as
one.

\paragraph{What the impact must not become.} \evint{} Three rules hold whatever the results.
\begin{itemize}[nosep]
  \item An unvalidated gap candidate never reaches a learner (\cref{sec:education}).
  \item Organizational outputs are aggregate by default, never feed performance evaluation, and name
    a person only with consent (\cref{sec:organizations}).
  \item Coverage is reported as a diagnostic, never a key performance indicator. As a target, it
    would reward documentation that raises the score over documentation that transfers knowledge.
\end{itemize}
